\documentclass[8pt,letterpaper]{article}
% ============================================================
%  Math Camp Formula Sheet -- TWO PAGES (US Letter, front & back)
%  Compile: pdflatex cheatsheet2.tex
%  DENSITY KNOB: \bodysize below. Smaller number = more fits.
%    6.6 -> roomy | 6.2 -> current | 5.8 -> very tight
% ============================================================
\usepackage[utf8]{inputenc}
\usepackage[T1]{fontenc}
\usepackage[letterpaper,margin=0.30in]{geometry}
\usepackage{amsmath,amssymb,amsfonts}
\usepackage{multicol}
\usepackage{enumitem}
\usepackage{xcolor}
\usepackage{array}
\usepackage{mathtools}

\makeatletter
\pagestyle{empty}
\clubpenalty=10000
\widowpenalty=10000
\displaywidowpenalty=10000
\setlength{\parindent}{0pt}
\setlength{\parskip}{0pt}
\setlength{\columnsep}{9pt}
\setlength{\columnseprule}{0.3pt}
\renewcommand{\columnseprulecolor}{\color{gray}}
\setlength{\abovedisplayskip}{1.5pt}
\setlength{\belowdisplayskip}{1.5pt}
\setlength{\abovedisplayshortskip}{1pt}
\setlength{\belowdisplayshortskip}{1pt}

\newcommand{\bodysize}{\fontsize{6.15}{6.95}\selectfont}

\newcommand{\sect}[1]{%
  \vspace{2.2pt}%
  \noindent\colorbox{black}{\parbox{\dimexpr\linewidth-2\fboxsep}{%
    \color{white}\bfseries\scshape\fontsize{6.4}{7.4}\selectfont
    \hyphenpenalty=10000\exhyphenpenalty=10000\raggedright #1}}%
  \vspace{1.4pt}\par\nopagebreak\@afterheading}

\newcommand{\keybox}[2]{%
  \vspace{1.5pt}%
  \noindent\fbox{\parbox{\dimexpr\linewidth-2\fboxsep-2\fboxrule}{%
    \bodysize\textbf{\textsc{#1}}\ #2}}%
  \vspace{1.5pt}\par}

\newenvironment{tl}
  {\begin{itemize}[leftmargin=6pt,itemsep=0.3pt,topsep=1pt,parsep=0pt,
                   partopsep=0pt,label=\textcolor{gray}{\tiny$\bullet$}]}
  {\end{itemize}}

\newcommand{\R}{\mathbb{R}}
\newcommand{\Rn}{\mathbb{R}^n}
\newcommand{\F}{\mathbb{F}}
\newcommand{\E}{\mathbb{E}}
\newcommand{\PP}{\mathbb{P}}
\newcommand{\Var}{\operatorname{Var}}
\newcommand{\Cov}{\operatorname{Cov}}
\newcommand{\tr}{\operatorname{tr}}
\newcommand{\rank}{\operatorname{rank}}
\newcommand{\Ker}{\operatorname{Ker}}
\newcommand{\Ima}{\operatorname{Im}}
\newcommand{\spn}{\operatorname{span}}
\newcommand{\intr}{\operatorname{int}}
\newcommand{\Dd}{D}
\newcommand{\tp}{^{\!\top}}
\newcommand{\ip}[2]{\langle #1,#2\rangle}
\newcommand{\nm}[1]{\lVert #1\rVert}

\makeatother
\begin{document}
\bodysize

\begin{center}
  \vspace{-6pt}
  {\fontsize{10}{11}\selectfont\bfseries Math Camp --- Formula Sheet}\quad
  {\itshape\fontsize{6}{7}\selectfont Linear Algebra $\cdot$ Analysis $\cdot$ Calculus $\cdot$ Optimization $\cdot$ Probability}
\end{center}
\vspace{-4pt}
\hrule height 0.9pt
\vspace{3pt}

\begin{multicols}{3}
\bodysize

% ================= LINEAR ALGEBRA =================
\sect{1. Spaces, Span, Independence}
\begin{tl}
\item \textbf{Subspace} $W\subseteq V$: $0\in W$, closed under $+$ and scalar $\cdot$. Intersections are subspaces; unions usually not.
\item $\spn\{v_1\dots v_m\}=\{\sum a_iv_i\}$ = smallest subspace containing them.
\item \textbf{LI}: $\sum a_iv_i=0\Rightarrow$ all $a_i=0$. \textbf{LD} $\iff$ some $v_j\in\spn$ of the others.
\item \textbf{Basis} = LI + spanning $\iff$ every $v$ has a \emph{unique} representation. $\dim V=\#$basis.
\item Any LI list $\le$ any spanning list. In $\dim n$: $n$ LI $\Rightarrow$ basis; $n$ spanning $\Rightarrow$ basis.
\item $\dim(U+W)=\dim U+\dim W-\dim(U\cap W)$.
\item \textbf{Direct sum} $V=U\oplus W$. Equivalent characterizations: (i) $V=U+W$ \emph{and} $U\cap W=\{0\}$; (ii) every $v\in V$ has a \emph{unique} $v=u+w$; (iii) the only way to write $0=u+w$ is $u=w=0$; (iv) $V=U+W$ and $\dim V=\dim U+\dim W$.
\item Basis of $U$ $\cup$ basis of $W$ is a basis of $U\oplus W$. Standard instances: $P^2=P\Rightarrow V=\Ker P\oplus\Ima P$; $V=U\oplus U^\perp$ for any subspace $U$ (orthogonal decomposition).
\end{tl}

\sect{2. Linear Maps}
\begin{tl}
\item $T(\alpha u+\beta v)=\alpha Tu+\beta Tv$; $\Ker T=\{v:Tv=0\}$, $\Ima T=T(V)$.
\item \textbf{$T$ injective $\iff \Ker T=\{0\}$.} Surjective $\iff \Ima T=W$.
\item \textbf{Rank--nullity:} $\dim V=\dim\Ker T+\dim\Ima T$.
\item $T:V\to V$, $\dim V<\infty$: \textbf{injective $\iff$ surjective $\iff$ invertible} (fails in $\infty$ dim).
\item Trick: $T(z_1\dots z_m)=\sum z_iv_i$. $\{v_i\}$ spans $\iff T$ onto; $\{v_i\}$ LI $\iff T$ 1-1.
\item $P^2=P\Rightarrow V=\Ker P\oplus\Ima P$, $v=(v-Pv)+Pv$.
\end{tl}

\sect{3. Matrices, Rank, Systems}
\begin{tl}
\item $\rank A=\#$pivots$=\dim\operatorname{Col}A=$ row rank $=$ col rank; $\dim\operatorname{Nul}A=n-\rank A$.
\item $Ax=b$ consistent $\iff \rank A=\rank[A|b] \iff b\in\operatorname{Col}A$. Solutions $=x_p+\operatorname{Nul}A$.
\item Basis of $\operatorname{Col}A$ = \emph{original} pivot columns. Basis of $\operatorname{Nul}A$: set each free var $=1$ in turn.
\item \textbf{$\rank A=1\iff A=cd\tp$}, i.e.\ $A_{jk}=c_jd_k$.
\item $\rank(AB)\le\min(\rank A,\rank B)$; $\rank A=\rank A\tp=\rank(A\tp A)$.
\item $(AB)\tp=B\tp A\tp$; $(AB)^{-1}=B^{-1}A^{-1}$.
\end{tl}

\keybox{Invertible $A$ ($n\times n$) --- all equivalent}{%
$\det A\neq0$ $\cdot$ $\rank n$ $\cdot$ $Ax=0$ only $x=0$ $\cdot$ $Ax=b$ unique $\forall b$ $\cdot$ cols LI $\cdot$ cols span $\F^n$ $\cdot$ $0$ not an eigenvalue $\cdot$ $A\tp$ invertible $\cdot$ $n$ pivots.}

\textbf{\textsc{Determinant / trace.}}
\begin{tl}
\item $\det(AB)=\det A\det B$; $\det A\tp=\det A$; $\det A^{-1}=1/\det A$; $\det(cA)=c^n\det A$.
\item Row swap flips sign; scaling a row by $c$ scales $\det$ by $c$; adding a multiple of a row: no change.
\item \textbf{Triangular/diagonal:} $\det=\prod$ diag., eigenvalues $=$ diagonal entries.
\item $2\times2$: $\det=ad-bc$, $A^{-1}=\frac{1}{\det}\left(\begin{smallmatrix}d&-b\\-c&a\end{smallmatrix}\right)$.
\item $\tr A=\sum a_{ii}=\sum\lambda_i$; $\det A=\prod\lambda_i$; $\tr(AB)=\tr(BA)$.
\item \textbf{Symmetric} $A=A\tp$; \textbf{orthogonal} $Q\tp Q=I$ ($Q^{-1}=Q\tp$, $\det=\pm1$, preserves norms); \textbf{idempotent} $P^2=P$ ($\lambda\in\{0,1\}$).
\end{tl}

\sect{4. Eigenvalues \& Diagonalization}
\begin{tl}
\item $Av=\lambda v$, $v\neq0$. \textbf{Eigenvalue $\iff$ kernel:}
\begin{align*}
\lambda\text{ eigenvalue}\ &\iff\ \Ker(A-\lambda I)\neq\{0\}\\
&\iff\ A-\lambda I\ \text{singular}\\
&\iff\ \det(A-\lambda I)=0.
\end{align*}
So $E_\lambda=\Ker(A-\lambda I)$ and geo.\ mult.\ $=\dim\Ker(A-\lambda I)=n-\rank(A-\lambda I)$.
\item Special case $\lambda=0$: $0\in\sigma(A)\iff\Ker A\neq\{0\}\iff A$ \emph{not} invertible. Eigenvectors for $\lambda=0$ are exactly the nonzero elements of $\Ker A$.
\item $2\times2$: $\lambda^2-(\tr A)\lambda+\det A=0$.
\item Distinct $\lambda$ $\Rightarrow$ eigenvectors LI $\Rightarrow$ $n$ distinct $\lambda\Rightarrow$ diagonalizable (\emph{sufficient, not necessary}).
\item $1\le$ geo.\ mult.\ $(\dim E_\lambda)\le$ alg.\ mult. \textbf{Diagonalizable $\iff$ geo $=$ alg for every $\lambda$ $\iff$ $n$ LI eigenvectors.} Repeated $\lambda$ alone decides nothing ($I$ vs $\left(\begin{smallmatrix}0&1\\0&0\end{smallmatrix}\right)$).
\item $X=[v_1\cdots v_n]$ (eigenvectors in \emph{columns}), $D=\operatorname{diag}(\lambda_i)$ \emph{in the same order}. $AX=XD$, $A=XDX^{-1}$, \textbf{$A^k=XD^kX^{-1}$}.
\item \textbf{Similarity} $B=S^{-1}AS$: same $\lambda,\det,\tr,\rank$, char.\ poly. $Av=\lambda v\Rightarrow B(S^{-1}v)=\lambda(S^{-1}v)$.
\item \textbf{Spectral thm (real symmetric):} $\lambda$ all real; eigenvectors of distinct $\lambda$ are $\perp$; $A=Q\Lambda Q\tp$.
\end{tl}

\sect{5. Quadratic Forms \& Definiteness}
$Q(x)=x\tp Ax$, $A$ symmetric (else use $\tfrac12(A+A\tp)$).
\begin{center}\fontsize{5.9}{6.8}\selectfont
\begin{tabular}{@{}p{0.13\linewidth}p{0.28\linewidth}p{0.45\linewidth}@{}}
\hline
 & eigenvalues & minors\\ \hline
PD & all $\lambda>0$ & leading $D_k>0$\\
ND & all $\lambda<0$ & $D_1<0,D_2>0,\dots$\\
PSD & all $\lambda\ge0$ & \emph{all} principal minors $\ge0$\\
NSD & all $\lambda\le0$ & $(-1)^k\cdot$minor $\ge0$\\ \hline
\end{tabular}
\end{center}
\begin{tl}
\item \textbf{Sylvester's criterion} (definite case only): $A$ sym.\ is PD $\iff D_1>0,\dots,D_n>0$, where $D_k=\det$ of the top-left $k\times k$ block. ND: apply to $-A$, i.e.\ signs alternate $D_1<0,D_2>0,\dots$
\item \textbf{Leading} $D_k$: only the $n$ top-left blocks. \textbf{Principal} minor: pick \emph{any} index set, keep those rows \emph{and} the same columns ($2^n-1$ of them). Leading $\subset$ principal.
\item \colorbox{black!12}{Sylvester with $\ge$ is FALSE.} $\operatorname{diag}(0,-1)$: $D_1=D_2=0\ge0$ but it is NSD. The principal minor $\{2\}=-1$ catches it. So \textbf{PSD requires all principal minors $\ge0$}, not just leading.
\item $2\times2$ shortcut: PD $\iff a>0$ and $\det>0$; ND $\iff a<0$ and $\det>0$; indefinite $\iff\det<0$.
\item Eigenvalue test never has this trap --- use it when $\lambda$ are cheap. $A\tp A$ always PSD; PD $\iff A$ has LI columns.
\end{tl}
\keybox{After you classify $H$ --- what it buys you}{%
\emph{At a critical point $x^*$ (i.e.\ $\nabla f(x^*)=0$):} PD $\Rightarrow$ strict local \textbf{min}; ND $\Rightarrow$ strict local \textbf{max}; indefinite $\Rightarrow$ \textbf{saddle}; PSD/NSD but singular $\Rightarrow$ \textbf{inconclusive} (e.g.\ $x^4$ vs $-x^4$ vs $x^3$ at $0$).\\[1pt]
\emph{On all of $U$ (not just at $x^*$):} $H$ PSD everywhere $\Rightarrow f$ convex; $H$ ND/NSD everywhere $\Rightarrow f$ concave --- and then the local verdict upgrades to \textbf{global}.}

\sect{6. Inner Products, Gram--Schmidt, Projection}
\begin{tl}
\item \textbf{Euclidean inner product} on $\Rn$: $\ip{u}{v}=\sum_{i=1}^n u_iv_i=u\tp v$, and $\nm{v}=\sqrt{\ip{v}{v}}=\big(\sum_i v_i^2\big)^{1/2}$.
\item Axioms: linear in 1st arg, (conjugate-)symmetric, $\ip{v}{v}>0$ for $v\neq0$. Other instances: $\ip{p}{q}=\int_0^1pq\,dx$ on $P_2(\R)$; $\ip{X}{Y}=\E[XY]$ on random variables.
\item $u\perp v\iff\ip{u}{v}=0$. $U^\perp=\{v:\ip{v}{u}=0\ \forall u\in U\}$ is a subspace; $V=U\oplus U^\perp$.
\item \textbf{Cauchy--Schwarz} $|\ip{u}{v}|\le\nm{u}\nm{v}$ (= iff LD). \textbf{Triangle} $\nm{u+v}\le\nm{u}+\nm{v}$. \textbf{Reverse} $\big|\nm{u}-\nm{v}\big|\le\nm{u-v}$.
\item $u\perp v\Rightarrow\nm{u+v}^2=\nm{u}^2+\nm{v}^2$.
\item \textbf{Operator norm} $\nm{A}=\max_{\nm{v}=1}\nm{Av}$, so $\nm{Ah}\le\nm{A}\nm{h}$ (\emph{not} Cauchy--Schwarz: matrix $\times$ vector, not an inner product).
\item \textbf{Gram--Schmidt (full form --- subtract the projection onto each previous vector):}
\begin{align*}
u_1&=v_1,\\
u_k&=v_k-\sum_{j<k}\frac{\ip{v_k}{u_j}}{\ip{u_j}{u_j}}\,u_j,\\
e_k&=u_k/\nm{u_k}.
\end{align*}
Each term $\frac{\ip{v_k}{u_j}}{\ip{u_j}{u_j}}u_j$ is the projection of $v_k$ onto $u_j$; you strip off every old direction, and what remains is $\perp$ to all of them. Normalize \emph{last} (or normalize as you go, and the denominators become $1$). $\spn\{u_1..u_k\}=\spn\{v_1..v_k\}$ at every step.
\item \textbf{Projection} onto $U$, o.n.\ basis: $P_Ub=\sum_j\ip{b}{e_j}e_j$. Residual $b-P_Ub\perp U$; $P_Ub$ unique closest point of $U$ to $b$.
\item Least squares: $A\tp A\hat x=A\tp b$, $P=A(A\tp A)^{-1}A\tp$.
\item \textbf{$\ip{u}{v}=0\iff\nm{u}\le\nm{u+av}\ \forall a$} --- the least-squares FOC.
\item $P_2(\R)$: $\ip{p}{q}=\int_0^1pq\,dx$; $\int_0^1x^k=\frac{1}{k+1}$.
\end{tl}

% ================= TOPOLOGY =================
\sect{7. Norms, Sequences, sup/inf}
\begin{tl}
\item All norms on $\Rn$ are equivalent $\Rightarrow$ they give the same convergent sequences, open sets and continuous functions. Work with whichever is convenient.
\item \textbf{$x_k\to x$}: $\forall\varepsilon>0\ \exists N: k\ge N\Rightarrow\nm{x_k-x}<\varepsilon$. In $\Rn$ $\iff$ coordinatewise.
\item Limits unique; convergent $\Rightarrow$ bounded; every subsequence has the same limit.
\item \textbf{$x_k\not\to x$} (negation, used constantly): $\exists\varepsilon>0$ and a subsequence with $\nm{x_{k_j}-x}\ge\varepsilon\ \forall j$.
\item \textbf{Cauchy}: $\forall\varepsilon\ \exists N: m,k\ge N\Rightarrow\nm{x_m-x_k}<\varepsilon$. $\Rn$ complete: Cauchy $\iff$ convergent.
\item \textbf{Bolzano--Weierstrass:} every bounded sequence in $\Rn$ has a convergent subsequence.
\item Monotone $+$ bounded (in $\R$) $\Rightarrow$ convergent.
\item \textbf{$m=\inf S$}: (i) $m\le s\ \forall s\in S$; (ii) $m$ is the \emph{greatest} such. Working form: $\forall\varepsilon>0\ \exists s\in S$, $s<m+\varepsilon$ (take $\varepsilon=1/k$ to get $s_k\to m$). Sup mirrors. Inf need not be attained: $\inf(0,1]=0\notin(0,1]$.
\item Completeness axiom: nonempty $S\subseteq\R$ bounded below has $\inf S\in\R$. This is what underlies Bolzano--Weierstrass, monotone convergence and Heine--Borel.
\item \textbf{Axiom of choice:} given any family $\{A_i\}_{i\in I}$ of nonempty sets, there is a choice function $c$ with $c(i)\in A_i$ for every $i$ --- i.e.\ you may simultaneously ``pick one element from each'' even with no rule for picking. Every proof that says \emph{choose $x_k\in B_{1/k}(x)\cap A$ for each $k$} (Pb.\ 3) or \emph{choose $x_k$ with $f(x_k)=y_k$} (Pb.\ 5) uses the countable version. Equivalent to Zorn's lemma and to the well-ordering theorem; needed to prove every vector space has a basis.
\end{tl}

\sect{8. Open, Closed, Closure}
\begin{tl}
\item $B_\varepsilon(x)=\{y:\nm{y-x}<\varepsilon\}$. \textbf{$A$ open} $\iff\forall x\in A\ \exists\varepsilon>0$, $B_\varepsilon(x)\subseteq A$.
\item \textbf{$A$ closed} $\iff A^c$ open $\iff$ ($x_k\in A$, $x_k\to x\Rightarrow x\in A$) $\iff$ \textbf{$A=\bar A$} $\iff$ \textbf{$A$ contains all of its limit points} ($A'\subseteq A$).
\item \textbf{Interior} $\intr A$ = largest open set $\subseteq A$ = $\{x: \exists\varepsilon, B_\varepsilon(x)\subseteq A\}$. \textbf{$A$ open $\iff A=\intr A$.}
\item \textbf{Closure} $\bar A$ = smallest closed set $\supseteq A$ = $A\cup A'$ = $\{x: \forall\varepsilon,\ B_\varepsilon(x)\cap A\ne\emptyset\}$ = all limits of sequences in $A$.
\item \textbf{Boundary} $\partial A=\bar A\setminus\intr A=\{x:\forall\varepsilon,\ B_\varepsilon(x)$ meets both $A$ and $A^c\}$. Always closed. $\bar A=A\cup\partial A$.
\item \textbf{Limit (accumulation) point} $x\in A'$: every $B_\varepsilon(x)$ contains a point of $A$ other than $x$. \textbf{Isolated}: some $B_\varepsilon(x)\cap A=\{x\}$.
\item $\overline{A^c}=(\intr A)^c$; $\intr(A^c)=(\bar A)^c$ (De Morgan for topology).
\item Arbitrary $\bigcup$ of open is open; finite $\bigcap$ of open is open (dually for closed). $\emptyset,\Rn$ are both open and closed.
\item \colorbox{black!12}{Infinite $\bigcap$ of open need not be open:} $\bigcap_k(-\frac1k,\frac1k)=\{0\}$.
\item \textbf{Bounded}: $\exists M$ with $\nm{x}\le M\ \forall x\in A$. \textbf{Dense}: $\bar A=\Rn$ (e.g.\ $\mathbb Q$ in $\R$).
\end{tl}

\sect{9. Compactness}
\begin{tl}
\item \textbf{Compact} (Heine--Borel in $\Rn$): closed \emph{and} bounded $\iff$ every open cover has a finite subcover $\iff$ every sequence has a subsequence converging \emph{in} the set.
\item Closed subset of compact is compact (bounded from $K$; limit retained by closedness of $F$).
\item Nested nonempty compacts $K_1\supseteq K_2\supseteq\cdots\Rightarrow\bigcap K_k\neq\emptyset$.
\item Compact $\Rightarrow$ $\sup$ and $\inf$ are attained (belong to the set).
\end{tl}

\sect{10. Limits \& Continuity}
\begin{tl}
\item $\lim_{x\to a}f(x)=L$: $\forall\varepsilon>0\ \exists\delta>0: 0<\nm{x-a}<\delta\Rightarrow\nm{f(x)-L}<\varepsilon$.
\item \textbf{$f$ continuous at $a$}: $\forall\varepsilon\ \exists\delta:\nm{x-a}<\delta\Rightarrow\nm{f(x)-f(a)}<\varepsilon$ $\iff$ $x_k\to a\Rightarrow f(x_k)\to f(a)$.
\item \textbf{Global:} $f$ continuous $\iff f^{-1}(\text{open})$ open $\iff f^{-1}(\text{closed})$ closed. \colorbox{black!12}{Preimages, not images.}
\item Corollaries: $\{f>c\},\{f<c\}$ open; $\{f\ge c\},\{f\le c\},\{f=c\}$ closed.
\item Sums, products, quotients (nonzero denom.), compositions of continuous are continuous.
\item \textbf{Weierstrass (EVT):} $K$ compact, $f$ cont.\ $\Rightarrow f(K)$ compact $\Rightarrow$ ($m=1$) max and min attained.
\item \textbf{IVT} ($n=m=1$): $f$ cont.\ on $[a,b]$, $y$ between $f(a),f(b)$ $\Rightarrow\exists c$ with $f(c)=y$.
\item Continuous bijection on compact $K$ $\Rightarrow f^{-1}$ continuous.
\item \textbf{Uniform continuity}: $\delta$ independent of $x$. Continuous on compact $\Rightarrow$ uniformly continuous. \textbf{Lipschitz}: $\nm{f(x)-f(y)}\le L\nm{x-y}$ $\Rightarrow$ unif.\ cont.
\item $d(x,A)=\inf_{a\in A}\nm{x-a}$ is 1-Lipschitz; $d(x,A)=0\iff x\in\bar A$.
\end{tl}

% ================= ONE-VARIABLE CALCULUS =================
\sect{11. One-Variable Calculus}
\begin{tl}
\item $f'(a)=\lim_{h\to0}\frac{f(a+h)-f(a)}{h}$. Differentiable $\Rightarrow$ continuous (converse false: $|x|$ at $0$).
\item \textbf{Product:} $(fg)'=f'g+fg'$. \textbf{Quotient:} $\left(\frac fg\right)'=\frac{f'g-fg'}{g^2}$. \textbf{Chain:} $(f\circ g)'(x)=f'(g(x))g'(x)$.
\item \textbf{Inverse:} $(f^{-1})'(y)=\dfrac{1}{f'(f^{-1}(y))}$, valid when $f'\ne0$ there.
\end{tl}
\begin{center}\fontsize{5.9}{6.9}\selectfont
\begin{tabular}{@{}p{0.30\linewidth}p{0.30\linewidth}@{}}
\hline
$f$ & $f'$\\ \hline
$x^a$ & $ax^{a-1}$\\
$e^x$ & $e^x$\\
$a^x$ & $a^x\ln a$\\
$\ln x$ & $1/x$\\
$\log_a x$ & $1/(x\ln a)$\\
$\sin x$ & $\cos x$\\
$\cos x$ & $-\sin x$\\
\hline
\end{tabular}
\end{center}
\begin{tl}
\item \textbf{Rolle:} $f$ cont.\ $[a,b]$, diff.\ $(a,b)$, $f(a)=f(b)$ $\Rightarrow\exists c\in(a,b)$, $f'(c)=0$.
\item \textbf{MVT:} same hypotheses $\Rightarrow\exists c\in(a,b)$ with $f'(c)=\dfrac{f(b)-f(a)}{b-a}$.
\item Consequences: $f'\equiv0\Rightarrow f$ const; $f'>0\Rightarrow$ strictly increasing; $|f'|\le L\Rightarrow$ $L$-Lipschitz.
\end{tl}

% ================= MULTIVARIABLE =================
\sect{12. Differentiation in $\Rn$}
\begin{tl}
\item \textbf{Definition:} $f(x+h)=f(x)+\Dd f(x)h+r(h)$, $\nm{r(h)}/\nm{h}\to0$. \textbf{$\Dd f$ unique} (put $h=tv$, divide by $t$, $t\to0$; then $(A-B)v=0\ \forall v$).
\item Read it as: \emph{$f$ is affine to first order near $x$}; $\Dd f(x)h$ is the predicted change, $r(h)$ the leftover.
\end{tl}
\keybox{Jacobian ($m\times n$)}{%
$\Dd f(x)=\begin{pmatrix}
\partial_1f_1 & \cdots & \partial_nf_1\\
\vdots & & \vdots\\
\partial_1f_m & \cdots & \partial_nf_m\end{pmatrix}$,
entry $(i,j)=\dfrac{\partial f_i}{\partial x_j}$. \textbf{Row $i$ = $\nabla f_i\tp$}; column $j$ = derivatives w.r.t.\ $x_j$. $m=1$: $\Dd f=\nabla f\tp$ (a row).}
\keybox{Hessian ($n\times n$)}{%
$H_f(x)=\Dd(\nabla f)(x)=\begin{pmatrix}
f_{11} & \cdots & f_{1n}\\
\vdots & & \vdots\\
f_{n1} & \cdots & f_{nn}\end{pmatrix}$,
$f_{ij}=\dfrac{\partial^2f}{\partial x_i\partial x_j}$. \textbf{Symmetric if $f\in C^2$} (Young/Clairaut).}
\begin{tl}
\item \textbf{Directional:} $D_vf(x)=\lim_{t\to0}\frac{f(x+tv)-f(x)}{t}$. If differentiable, $D_vf=\Dd f(x)v$ (\emph{linear in $v$}). Partials $=D_{e_j}f$.
\item Differentiable $\Rightarrow$ continuous. Partials exist $\nRightarrow$ continuous.
\item \textbf{$C^1\Rightarrow$ differentiable}: split the increment one variable at a time, MVT on each piece $\Rightarrow f_x(\xi,b+k)h+f_y(a,\eta)k$; subtract $f_x(a,b)h+f_y(a,b)k$; bound by $(|\Delta f_x|+|\Delta f_y|)\sqrt{h^2+k^2}$; continuity kills each $\Delta$.
\item \textbf{Chain rule:} $\Dd(g\circ f)(x)=\Dd g(f(x))\,\Dd f(x)$, sizes $(p\times m)(m\times n)=p\times n$. Order matters; outer evaluated \emph{at} $f(x)$. Inner Jacobian has \textbf{one column per independent variable}.
\item Scalar form: $\dfrac{\partial (g\circ f)}{\partial x_j}=\sum_i\dfrac{\partial g}{\partial y_i}\dfrac{\partial f_i}{\partial x_j}$. If $x$ enters twice, both paths contribute (total vs partial derivative).
\item $\nabla(a\tp x)=a$; $\Dd(Ax)=A$; $\nabla(x\tp Ax)=(A+A\tp)x=2Ax$ if $A$ sym.; $H(x\tp Ax)=2A$.
\item $\nabla f$ points in the direction of steepest ascent; $\nabla f\perp$ level set $\{f=c\}$.
\end{tl}

\sect{13. Restrict to a Line}
$g(t)=f(a+t(b-a))$, $t\in[0,1]$: \emph{multivariate $\to$ segment $\to$ one-variable calculus.}
\begin{tl}
\item $g'(t)=\Dd f(a+t(b-a))(b-a)$.
\item \textbf{MVT:} $f$ real-valued $\Rightarrow\exists c$ on the segment with $f(b)-f(a)=\Dd f(c)(b-a)$. \colorbox{black!12}{Fails for $m\ge2$} (each component gets its own $c$); use $|f(b)-f(a)|\le\sup_{z\in[a,b]}\nm{\Dd f(z)}\nm{b-a}$.
\item $U$ convex, $\Dd f\equiv0\Rightarrow f$ constant. $H_f\equiv0\Rightarrow\Dd(\nabla f)\equiv0\Rightarrow\nabla f\equiv a\Rightarrow f(x)=a\tp x+b$ (apply the previous line to $g(x)=f(x)-a\tp x$).
\item \textbf{Second-order Taylor approximation} (the working form): for $h$ small,
\[ f(x+h)\;\approx\; f(x)+\nabla f(x)\tp h+\tfrac12\,h\tp H_f(x)\,h, \]
with error $o(\nm{h}^2)$. First-order: $f(x+h)\approx f(x)+\nabla f(x)\tp h$ (the tangent plane).
\item Everything is evaluated \emph{at $x$}: gradient $\nabla f(x)$ (a vector, paired with $h$) and Hessian $H_f(x)$ (a matrix, sandwiched as $h\tp H h$). The second-order term is a \emph{scalar quadratic form} because $f$ is scalar-valued.
\item At a critical point $\nabla f(x^*)=0$, so $f(x^*+h)-f(x^*)\approx\frac12h\tp H_fh$: the sign of the quadratic form \emph{is} the local min/max verdict. That is why definiteness is the SOC.
\end{tl}

\sect{14. Inverse \& Implicit Function Theorems}
\keybox{Inverse FT}{$f:\Rn\to\Rn$ $C^1$ near $x_0$, $\det \Dd f(x_0)\neq0$ $\Rightarrow$ $\exists$ open $U\ni x_0$, $V\ni f(x_0)$ with $f:U\to V$ a $C^1$ bijection and
\[\Dd(f^{-1})(f(x_0))=[\Dd f(x_0)]^{-1}.\]
\textbf{Local only; sufficient, not necessary.}}
\begin{tl}
\item $F(x,y)=(e^x\cos y,e^x\sin y)$: $\det \Dd F=e^{2x}\neq0$ everywhere, yet $F(x,y+2\pi)=F(x,y)$ --- locally 1-1, not globally. IFT is local.
\item $f(x)=x^3$: $f'(0)=0$ but globally 1-1; $f^{-1}=x^{1/3}$ \emph{not} differentiable at $0$. Singular $\Dd f$ rules out a \emph{differentiable} inverse, not injectivity.
\end{tl}
\keybox{Implicit FT}{$F:\R^n\times\R^m\to\R^m$, $F(x_0,y_0)=0$, $F\in C^1$, $\Dd_yF(x_0,y_0)$ \emph{invertible} $\Rightarrow$ locally a unique $C^1$ map $y=g(x)$ with $g(x_0)=y_0$ and
\[\Dd g(x)=-[\Dd_yF]^{-1}\Dd_xF.\]
Scalar case $F(x,y,z)=0$, $F_z\ne0$: $g_x=-F_x/F_z$, $g_y=-F_y/F_z$.}
\begin{tl}
\item Sphere $x^2+y^2+z^2=1$: at $p=(0,0,1)$, $F_z=2z\neq0\Rightarrow z=g(x,y)$; $F_x=0\Rightarrow$ cannot solve $x=h(y,z)$. At $q=(1,0,0)$ it reverses (tangent plane vertical in that variable).
\item Derive by chain rule, not memory: $F(x,y,g(x,y))=0\Rightarrow F_x+F_zg_x=0$.
\item $GL_n=\det^{-1}(\R\setminus\{0\})$ is open $\Rightarrow$ $\Dd f(x^*)$ invertible $\Rightarrow \Dd f(x)$ invertible nearby.
\end{tl}

% ================= OPTIMIZATION =================
\sect{15. Unconstrained Optimization}
\begin{tl}
\item \textbf{FOC:} interior local extremum $+$ differentiable $\Rightarrow\nabla f(x^*)=0$ (necessary, not sufficient: $x^3$ at $0$).
\item \textbf{SOC (local, at a critical point $x^*$):} evaluate $H_f(x^*)$ and classify --- PD $\Rightarrow$ strict local min; ND $\Rightarrow$ strict local max; indefinite $\Rightarrow$ saddle; singular semidefinite $\Rightarrow$ inconclusive.
\end{tl}
\keybox{Convex vs.\ concave --- the Hessian test}{%
Let $U$ be convex and $f\in C^2$ on $U$. Then, \emph{for all $x\in U$}:
\begin{center}\fontsize{5.9}{6.9}\selectfont
\begin{tabular}{@{}p{0.26\linewidth}p{0.24\linewidth}p{0.34\linewidth}@{}}
\hline
$H_f(x)$ & $f$ is & optimum \\ \hline
PSD $\forall x$ & convex & crit.\ pt.\ = global \textbf{min}\\
PD $\forall x$ & strictly convex & unique global min\\
NSD $\forall x$ & concave & crit.\ pt.\ = global \textbf{max}\\
ND $\forall x$ & strictly concave & unique global max\\ \hline
\end{tabular}
\end{center}
\vspace{1pt}
\textbf{Quantifier matters:} PD \emph{at one point $x^*$} gives only a \emph{local} verdict; PSD \emph{everywhere on $U$} gives convexity, and then $\nabla f(x^*)=0\Rightarrow x^*$ is a \textbf{global} min.}
\begin{tl}
\item Equivalent definitions of convex (any one suffices): $f(\lambda x+(1-\lambda)y)\le\lambda f(x)+(1-\lambda)f(y)$ for $\lambda\in[0,1]$ (chord above graph); $f(y)\ge f(x)+\nabla f(x)\tp(y-x)$ (graph above tangent); $H_f$ PSD. Concave: reverse every inequality, or apply these to $-f$.
\item Convex $f$: the set of minimizers is convex; every local min is global; there are no saddles and no strict local maxima in the interior.
\item Existence: $f$ continuous on compact $\Rightarrow$ max and min attained (Weierstrass).
\end{tl}

\sect{16. Lagrange \& KKT}
\textbf{Equality constraints.} $\max f(x)$ s.t.\ $h_j(x)=0$: $\mathcal L=f-\sum\lambda_jh_j$; FOC $\nabla f=\sum\lambda_j\nabla h_j$, $h_j=0$; $\lambda$ free in sign; CQ: $\{\nabla h_j\}$ LI at $x^*$.
\keybox{KKT --- $\max f(x)$ s.t.\ $g_i(x)\ge0$}{$\mathcal L=f(x)+\sum_i\lambda_ig_i(x)$\\[1pt]
(1) \textbf{Stationarity} $\nabla f(x^*)+\sum_i\lambda_i\nabla g_i(x^*)=0$\\
(2) \textbf{Primal feas.} $g_i(x^*)\ge0$\\
(3) \textbf{Dual feas.} $\lambda_i\ge0$\\
(4) \textbf{Compl.\ slackness} $\lambda_ig_i(x^*)=0$\\[1pt]
$\lambda_i>0\Rightarrow$ constraint binds; slack $\Rightarrow\lambda_i=0$.}
\begin{tl}
\item For a \textbf{min}: minimize $-f$, or use $\mathcal L=f-\sum\lambda_ig_i$ with the same signs on $\lambda,g$.
\item Rewrite every constraint as $g\ge0$ first (e.g.\ $w-px\ge0$, $x\ge0$) so the conventions hold.
\item Case method: guess which constraints bind, solve, then \emph{check} $\lambda\ge0$ and $g\ge0$.
\item CQ (Slater / LI of active gradients) makes KKT \emph{necessary}. $f$ concave, $g_i$ concave $\Rightarrow$ KKT also \textbf{sufficient}, global max.
\item \textbf{Multiplier:} $\lambda_i=\partial V/\partial c_i$ = shadow price. \textbf{Envelope:} $dV/d\theta=\partial\mathcal L/\partial\theta$ at $(x^*,\lambda^*)$.
\item \textbf{Bordered Hessian} ($n$ vars, $m$ constraints): check the last $n-m$ leading principal minors of $\bar H=\left(\begin{smallmatrix}0&\Dd h\\ \Dd h\tp&H_{\mathcal L}\end{smallmatrix}\right)$. Max: signs alternate starting $(-1)^{m+1}$. Min: all of sign $(-1)^m$.
\end{tl}

% ================= PROBABILITY =================
\sect{17. Probability Basics}
\begin{tl}
\item Sample space $\Omega$, events $\mathcal F$, measure $\PP$: $\PP(A)\ge0$, $\PP(\Omega)=1$, countable additivity on disjoint sets.
\item $\PP(A^c)=1-\PP(A)$; $\PP(A\cup B)=\PP(A)+\PP(B)-\PP(A\cap B)$; monotone: $A\subseteq B\Rightarrow\PP(A)\le\PP(B)$.
\item \textbf{Conditional:} $\PP(A\mid B)=\dfrac{\PP(A\cap B)}{\PP(B)}$, $\PP(B)>0$. \textbf{Multiplication:} $\PP(A\cap B)=\PP(A\mid B)\PP(B)$.
\item \textbf{Independence:} $\PP(A\cap B)=\PP(A)\PP(B)\iff\PP(A\mid B)=\PP(A)$. Pairwise $\nRightarrow$ mutual.
\item \textbf{Total probability:} $\{B_i\}$ a partition $\Rightarrow\PP(A)=\sum_i\PP(A\mid B_i)\PP(B_i)$.
\item \textbf{Bayes:} $\PP(B_j\mid A)=\dfrac{\PP(A\mid B_j)\PP(B_j)}{\sum_i\PP(A\mid B_i)\PP(B_i)}$.
\item \textbf{CDF} $F(x)=\PP(X\le x)$: nondecreasing, right-continuous, $F(-\infty)=0$, $F(\infty)=1$. Continuous case $f=F'$, $\PP(a<X\le b)=F(b)-F(a)$.
\end{tl}

\sect{18. Expectation \& Moments}
\begin{tl}
\item $\E[X]=\sum_x x\,p(x)$ or $\int x f(x)dx$. \textbf{LOTUS:} $\E[g(X)]=\int g(x)f(x)dx$ --- no need for the density of $g(X)$.
\item \textbf{Linearity} $\E[aX+bY+c]=a\E X+b\E Y+c$ \emph{always} (no independence needed).
\item $\Var X=\E[(X-\E X)^2]=\E[X^2]-(\E X)^2\ge0$; $\Var(aX+b)=a^2\Var X$.
\item \textbf{Covariance:} $\Cov(X,Y)=\E\big[(X-\E X)(Y-\E Y)\big]=\E[XY]-\E X\,\E Y$ (expand the product and use linearity to pass between the two). $\Cov(X,X)=\Var X$; symmetric and bilinear: $\Cov(aX+b,cY+d)=ac\,\Cov(X,Y)$.
\item $\Var(X+Y)=\Var X+\Var Y+2\Cov(X,Y)$; $\Var\big(\sum a_iX_i\big)=\sum_i\sum_j a_ia_j\Cov(X_i,X_j)$.
\item Independent $\Rightarrow$ $\E[XY]=\E X\E Y$ $\Rightarrow\Cov=0$. \colorbox{black!12}{Converse false} ($X\sim N(0,1)$, $Y=X^2$).
\item $\rho=\dfrac{\Cov(X,Y)}{\sigma_X\sigma_Y}\in[-1,1]$ (Cauchy--Schwarz); $|\rho|=1\iff$ $Y$ affine in $X$ a.s.
\item \textbf{Markov:} $X\ge0$, $a>0$ $\Rightarrow\PP(X\ge a)\le\E X/a$. \textbf{Chebyshev:} $\PP(|X-\mu|\ge k\sigma)\le1/k^2$.
\item \textbf{MGF} $M_X(t)=\E[e^{tX}]$; $M^{(k)}(0)=\E[X^k]$; determines the distribution; independent $\Rightarrow M_{X+Y}=M_XM_Y$.
\end{tl}

\keybox{Law of Iterated Expectations (LIE)}{%
\[\E[X]=\E\big[\E[X\mid Y]\big]\]
More generally $\E[X\mid Z]=\E\big[\E[X\mid Y,Z]\mid Z\big]$ (tower property; the \emph{coarser} conditioning wins).\\[1pt]
Key facts: $\E[X\mid Y]$ is a \emph{random variable} (a function of $Y$); $\E[g(Y)X\mid Y]=g(Y)\E[X\mid Y]$ (\emph{taking out what is known}); $X\perp Y\Rightarrow\E[X\mid Y]=\E X$.\\[1pt]
\textbf{Law of total variance:} $\Var X=\E[\Var(X\mid Y)]+\Var(\E[X\mid Y])$.}
\begin{tl}
\item $\E[X\mid Y]$ is the best predictor of $X$ given $Y$ in mean square: it minimizes $\E[(X-g(Y))^2]$ over all $g$. Residual $X-\E[X\mid Y]$ is uncorrelated with any $g(Y)$ --- the projection/orthogonality logic of \S6.
\item Discrete form: $\E[X]=\sum_y\E[X\mid Y=y]\PP(Y=y)$.
\end{tl}

\sect{19. Normal, LLN, CLT}
\begin{tl}
\item $X\sim N(\mu,\sigma^2)$: $\E X=\mu$, $\Var X=\sigma^2$. Affine maps stay normal: $aX+b\sim N(a\mu+b,\,a^2\sigma^2)$; in particular $Z=\frac{X-\mu}{\sigma}\sim N(0,1)$.
\item Sums of \emph{independent} normals are normal: $X\pm Y\sim N(\mu_X\pm\mu_Y,\ \sigma_X^2+\sigma_Y^2)$ (variances always add).
\item \textbf{LLN:} $\bar X_n\to\mu$ (in probability / almost surely) for i.i.d.\ with finite mean.
\item \textbf{CLT:} i.i.d., finite variance $\Rightarrow\sqrt n(\bar X_n-\mu)\xrightarrow{d}N(0,\sigma^2)$, i.e.\ $\bar X_n\approx N(\mu,\sigma^2/n)$ for large $n$.
\end{tl}

\sect{20. Counterexample Bank}
\begin{tl}
\item \textbf{Partials exist, not differentiable/continuous:} $f=\frac{xy}{x^2+y^2}$, $f(0,0)=0$: $f_x=f_y=0$ at $0$ but $f\to\frac12$ along $y=x$. (Polar: $f=\frac12\sin2\theta$, constant on rays.)
\item \textbf{All directional derivatives, still not differentiable:} $f=\frac{x^2y}{x^2+y^2}$, $f(0,0)=0$; $D_vf(0,0)=\frac{a^2b}{a^2+b^2}$ is \emph{not linear} in $v$.
\item \textbf{Continuous image of closed need not be closed:} $\frac{1}{1+x^2}$ on $\R$ gives $(0,1]$.
\item \textbf{Continuous image of bounded need not be bounded:} $1/x$ on $(0,1)$.
\item \textbf{No max:} $x$ on $(0,1)$ (bounded, not closed); $x$ on $[0,\infty)$ (closed, not bounded); $1/x$ on $(0,1]$ with $f(0)=0$ (compact, discontinuous).
\item \textbf{Continuous bijection, discontinuous inverse:} $t\mapsto(\cos t,\sin t)$ on $[0,2\pi)$ --- domain not compact.
\item \textbf{$\Dd f$ singular but 1-1:} $x^3$. \textbf{$\Dd f$ invertible, not globally 1-1:} $e^x(\cos y,\sin y)$.
\item \textbf{Not diagonalizable:} $\left(\begin{smallmatrix}0&1\\0&0\end{smallmatrix}\right)$ over $\mathbb C^2$.
\item \textbf{Uncorrelated but dependent:} $X\sim N(0,1)$, $Y=X^2$.
\end{tl}

\sect{21. Diagnose the Statement}
\begin{center}\fontsize{5.9}{6.9}\selectfont
\begin{tabular}{@{}p{0.84\linewidth}@{}p{0.08\linewidth}@{}}
\hline
claim & \\ \hline
differentiable $\Rightarrow$ continuous & T\\
all partials exist $\Rightarrow$ differentiable & F\\
all directional derivatives exist $\Rightarrow$ diff. & F\\
partials near $x$ $+$ continuous at $x$ $\Rightarrow$ diff. & T\\
continuous image of compact is compact & T\\
continuous image of closed is closed & F\\
closed subset of compact is compact & T\\
$\Dd f(x)$ invertible $\Rightarrow$ globally 1-1 & F\\
$\Dd f(x)$ singular $\Rightarrow$ no \emph{diff.} local inverse & T\\
$\Dd f(x)$ singular $\Rightarrow$ not locally 1-1 & F\\
convex $U$, $\Dd f\equiv0\Rightarrow f$ constant & T\\
convex $U$, $H_f\equiv0\Rightarrow f$ affine & T\\ \hline
\end{tabular}
\end{center}

\sect{22. Proof Moves}
\begin{tl}
\item Neighborhood property fails $\Rightarrow$ \textbf{build a sequence} with $\varepsilon=1/k$, extract, converge.
\item Set open/closed $\Rightarrow$ write it as $f^{-1}(\text{open/closed})$ for continuous $f$. Never form an image.
\item \textbf{Preimage skeleton:} $a\in f^{-1}(V)$ means $f(a)\in V$; openness of $V$ gives $\varepsilon$; continuity gives $\delta$; then $x\in B_\delta(a)\Rightarrow f(x)\in B_\varepsilon(f(a))\subseteq V\Rightarrow x\in f^{-1}(V)$. Conclude $B_\delta(a)\subseteq f^{-1}(V)$.
\item $x_k\not\to x$ $\Rightarrow$ $\exists\varepsilon>0$ and a subsequence bounded away; compactness gives a \emph{further} convergent subsequence; contradict via uniqueness of limits $+$ injectivity.
\item Bounding an $\inf$ you cannot evaluate: fix an arbitrary element, derive the inequality for all of them, then take $\inf$ on the side that does not depend on it.
\item $A=B$ for linear maps $\Rightarrow$ show $(A-B)v=0$ for all $v$.
\item Multivariate statement $\Rightarrow$ restrict to a line, use MVT or 1-D Taylor.
\item $a=b$ $\Rightarrow$ prove $\le$ and $\ge$, or $|a-b|<\varepsilon$ for every $\varepsilon>0$.
\item Showing something is \emph{not} linear: exhibit $v_1,v_2$ with $T(v_1+v_2)\ne T(v_1)+T(v_2)$.
\end{tl}

\end{multicols}
\end{document}
